% Generated draft: requires mathematical and source review.
\documentclass[11pt]{article}
\usepackage[margin=25mm]{geometry}
\usepackage{fontspec}
\setmainfont{Latin Modern Roman}
\usepackage{amsmath,amssymb,mathtools,unicode-math}
\setmathfont{Latin Modern Math}
\usepackage{enumitem,needspace,xurl,hyperref}
\setlength{\parindent}{0pt}
\setlength{\parskip}{5pt}
\setlength{\emergencystretch}{3em}
\newcommand{\N}{\mathbb{N}}
\newcommand{\Z}{\mathbb{Z}}
\newcommand{\Q}{\mathbb{Q}}
\newcommand{\R}{\mathbb{R}}
\newcommand{\C}{\mathbb{C}}
\newcommand{\F}{\mathbb{F}}
\newcommand{\A}{\mathbb{A}}
\newcommand{\PP}{\mathbb{P}}
\newcommand{\E}{\mathbb{E}}
\newcommand{\eps}{\varepsilon}
\newcommand{\dd}{\,\mathrm{d}}
\begin{document}
\section*{UnsolvedMath ID: 701469}
\textbf{Source problem code:} catalogue entry
\subsection{Definitions and mathematical statement}
For a group \(G\), use the commutator convention \([x,y]=xyx^{-1}y^{-1}\), and let \([G,G]\) denote the subgroup generated by all such commutators. For \(g\in[G,G]\), define its commutator length by
\[
\operatorname{cl}_G(g)=\min\left\{m\in\mathbb Z_{\geq0}:g=\prod_{i=1}^{m}[x_i,y_i],\quad x_i,y_i\in G\right\}.
\]
The empty product is the identity, so \(\operatorname{cl}_G(1)=0\). Define the stable commutator length by
\[
\operatorname{scl}_G(g)=\lim_{n\to\infty}\frac{\operatorname{cl}_G(g^n)}{n},\qquad n\in\mathbb Z_{\geq1}.
\]
This limit exists by subadditivity of \(n\mapsto\operatorname{cl}_G(g^n)\).

Interpret a hyperbolic group in the standard word-hyperbolic sense: a finitely generated group admitting a finite symmetric generating set whose Cayley graph, with edges of length one and its path metric, has \(\delta\)-thin geodesic triangles for some \(\delta\geq0\). Here thinness means that each side lies within distance \(\delta\) of the union of the other two sides.

\textbf{Conjecture (D. Calegari; Kourovka Notebook Problem 18.40).} For every hyperbolic group \(G\) and every \(g\in[G,G]\),
\[
\operatorname{scl}_G(g)\in\mathbb Q_{\geq0}.
\]
The supplied record labels this conjecture open; that label is supplied registry information, not an independently verified assessment.
\subsection{Short English statement}
D. Calegari's conjecture asks whether every element of the commutator subgroup of every word-hyperbolic group has rational stable commutator length. Stable commutator length is the limiting number of commutators needed to express powers of the element, divided by the exponent.
\subsection{Sources}
\begin{itemize}
\item Local UnsolvedMath registry, supplied complete record: permanent ID \(701469\), exact source code \texttt{catalogue entry}, title \emph{Kourovka Notebook Problem 18.40}. The record was created and updated on \texttt{2026-10-16T14:00:00Z}; its supplied status is \texttt{open}. These are registry metadata, not independent literature verification.
\item \emph{The Kourovka Notebook (Unsolved Problems in Group Theory)}, original Problem 18.40, proposed by D. Calegari in 2014. The supplied background identifies the 18th issue (2014) and the 21st issue edition (2026), \texttt{arXiv:1401.0300v48}. The supplied collection description attributes editing of Issues 1--21 (1965--2026) to E. I. Khukhro and V. D. Mazurov. This bibliographic attribution is reproduced from the record and has not been independently checked.
\end{itemize}
\end{document}
